
%%==================================================================
\section{The lower bound}\label{sec:lower}
%%==================================================================

The argument is that of \cite[Appendix C.1]{OAI26}, itself an instance of
Williams' method \cite{Wil13}; what changes is the circuit class, which is
handled by Theorem~\ref{thm:counting} in place of the $3$-CNF algorithm.

\subsection{Ingredients}
All machines are multitape Turing machines. We use three known results.

\begin{lemma}[Nondeterministic hierarchy, {\cite{Zak83}}; see {\cite[Theorem 2.1]{Wil13}}]\label{lem:zak}
There is a unary language $H\subseteq\{1\}^*$ with
$H\in\NTIME(2^n)\setminus\NTIME(2^{n/2})$.
\end{lemma}

\begin{lemma}[Local reduction, {\cite[Theorem 1 and Section 6]{JMV15}}]\label{lem:jmv}
Let $M$ be the machine deciding $H$ in time $2^n$. There are $q=n+O(\log n)$,
$v=n+O(\log n)$, a fixed integer $w$, and a deterministic polynomial-time
algorithm that, given $n$, outputs a map
\[
D_n:\{0,1\}^q\to\bigl(\{0,1\}^v\bigr)^3\times\{0,1\}^3
\]
with the following properties. Let $\varphi_n$ be the $3$-CNF on the variables
$\{X_\lambda:\lambda\in\{0,1\}^v\}$ whose clauses are, for $i\in\{0,1\}^q$ with
$D_n(i)=(\lambda_1,\lambda_2,\lambda_3,\epsilon_1,\epsilon_2,\epsilon_3)$, the
clauses $X_{\lambda_1}^{\epsilon_1}\vee X_{\lambda_2}^{\epsilon_2}\vee
X_{\lambda_3}^{\epsilon_3}$, where $X^1=X$ and $X^0=\neg X$. Then $\varphi_n$ is
satisfiable if and only if $1^n\in H$, and each output bit of $D_n$ depends on at
most $w$ bits of $i$.
\end{lemma}

The reduction of \cite{JMV15} applies to any nondeterministic computation of
length $T$, here $T=2^nn^{O(1)}$ including the overhead of the simulation, and
gives $\log T+O(\log\log T)$ index bits and clause bits computable by decision
trees of constant depth.

\begin{lemma}[Decoder circuits]\label{lem:decoder}
There is a circuit $\Delta_n$ that computes $D_n$, consists of $O(n)$ table
gates of fan-in at most $w$, and is constructed in time polynomial in $n$.
\end{lemma}

\begin{proof}
This is the argument of \cite[Appendix C.1]{OAI26}. For an output bit $o$ of
$D_n$, evaluate $D_n$ at the zero string and at all strings of Hamming weight
at most $w$; this is polynomially many evaluations since $w$ is fixed. If $o$
depends on coordinate $c$, there are two strings that differ only in $c$ and
give different values of $o$; setting to zero all coordinates outside the
support of $o$ keeps the values, so two such strings of weight at most $w$
exist and are among those evaluated. Hence the evaluations recover the support
of $o$ and its truth table on it. Each output bit is then one table gate of
fan-in at most $w$, and $D_n$ has $O(n)$ output bits.
\end{proof}

\subsection{The language $\Wl$}
Fix an integer $d\ge1$ with $v\le dn$ for all large $n$; $d=2$ will do. For a
label $\lambda\in\{0,1\}^v$ write $\bar\lambda\in\{0,1\}^{dn}$ for $\lambda$
followed by $dn-v$ zeros, and regard $\varphi_n$ as a formula on all $2^{dn}$
variables $X_{\bar\lambda}$, most of which do not occur. For $|i|=dn$, define
$\Wl(i)$ as bit $i$ of the lexicographically first satisfying assignment of
$\varphi_n$, with variables ordered by their labels, when $\varphi_n$ is
satisfiable, and $\Wl(i)=0$ otherwise; for lengths not divisible by $d$ put
$\Wl(i)=0$.

\begin{lemma}\label{lem:WinENP}
$\Wl\in\ENP$.
\end{lemma}

\begin{proof}
On an input of length $N=dn$, write out the $2^q$ clauses of $\varphi_n$ by
evaluating $D_n$, in time $2^{O(n)}$. Ask the oracle whether $\varphi_n$ is
satisfiable. If it is, determine the first assignment bit by bit: having fixed
a prefix, ask whether $\varphi_n$ with the prefix followed by zero has a
satisfying extension, and append zero or one accordingly. The queries have
length $2^{O(n)}$, and at most $2^{dn}$ of them determine bit $i$. The time is
$2^{O(N)}$.
\end{proof}

\subsection{From counting to wires}
\begin{theorem}[Transfer; Theorem D]\label{thm:transfer}
Let $\gamma>0$ be rational, let $A_1,b_1>0$, and suppose that a deterministic multitape Turing
machine decides, for every circuit $C\in\cK(s^{4/5},s^{2/3})$ with $s=s(C)\ge2$
given by a description of length $L$, whether $C$ is satisfiable, in time
$(s+L)^{b_1}2^{A_1s/(\log s)^\gamma}$. Then there is $c>0$ such that for infinitely
many $N$, every circuit in $\cK(N^{3/4},N^{1/2})$ with $N$ inputs deciding $\Wl$ on
$\{0,1\}^N$ has more than $cN(\log N)^\gamma$ wires.
\end{theorem}

\begin{proof}
Fix a rational $\eta$ with $0<\eta<1/(12A_1)$ and put
$s_n=\lfloor\eta n\lfloor\lfloor\log n\rfloor^\gamma\rfloor\rfloor$, which is
computable in time polynomial in $\log n$. Suppose,
for contradiction, that for all large $n$ some circuit $E_n\in\cK((dn)^{3/4},
(dn)^{1/2})$ with $dn$ inputs and at most $s_n$ wires decides $\Wl$ on
$\{0,1\}^{dn}$. We decide $H$ as follows. Reject inputs that are not unary. On
$1^n$, guess a circuit $E$ in this class with at most $s_n$ wires; its
description has length $L_E=2^{O(n^{3/4})}$, since it has at most $s_n$ gates,
each with a table of at most $2^{(dn)^{3/4}}$ bits or $O(s_n)$ numbers of
$O(\sqrt n)$ bits. Build the test circuit $\mathcal T$ with inputs
$i\in\{0,1\}^q$:
\begin{enumerate}[label=\textup{(\arabic*)},leftmargin=2em]
\item the decoder $\Delta_n$ of Lemma~\ref{lem:decoder} computes
$(\lambda_1,\lambda_2,\lambda_3,\epsilon_1,\epsilon_2,\epsilon_3)$;
\item three copies of $E$ read $\bar\lambda_1,\bar\lambda_2,\bar\lambda_3$, the
padding zeros coming from one constant gate;
\item one table gate of fan-in six outputs one if and only if every literal
$E(\bar\lambda_t)^{\epsilon_t}$ is false, that is, if the clause of index $i$ is
violated by the assignment $\lambda\mapsto E(\bar\lambda)$.
\end{enumerate}
Accept if and only if $\mathcal T$ is unsatisfiable.

If $1^n\in H$, then $\varphi_n$ is satisfiable and the guess $E=E_n$ computes a
satisfying assignment, so no clause is violated and the machine accepts. If
$1^n\notin H$, then every $E$ defines an assignment that violates some clause of
$\varphi_n$, so $\mathcal T$ is satisfiable and the machine rejects. Thus the
machine decides $H$.

The circuit $\mathcal T$ has $q=n+O(\log n)$ inputs and at most
$3s_n+O(n)$ wires, so $n\le s(\mathcal T)\le3s_n+B n$ for a constant $B$. Its
table gates have fan-in at most $\max(w,6,(dn)^{3/4})$ and its weights at most
$2^{(dn)^{1/2}}$. For $n\ge d^{15}$ we have $(dn)^{3/4}\le n^{4/5}$, and for
$n\ge d^3$ we have $(dn)^{1/2}\le n^{2/3}$; since $s(\mathcal T)\ge n$, the
circuit lies in $\cK(s^{4/5},s^{2/3})$ with $s=s(\mathcal T)$ for all large $n$.
By hypothesis its satisfiability is decided in time
\[
(s+L)^{b_1}2^{A_1s/(\log s)^\gamma}
\le2^{O(n^{3/4})}\cdot2^{A_1(3s_n+Bn)/(\log n)^\gamma}
=2^{(3A_1\eta+o(1))n},
\]
and $3A_1\eta<1/4$. Guessing $E$ and building $\mathcal T$ take
$2^{O(n^{3/4})}$ steps. Hence $H\in\NTIME(2^{n/2})$ after hardwiring finitely
many lengths, which contradicts Lemma~\ref{lem:zak}.

Therefore, for infinitely many $n$, every circuit in
$\cK((dn)^{3/4},(dn)^{1/2})$ deciding $\Wl$ on length $N=dn$ has more than $s_n$
wires. Since $\log n=\log N-\log d$, we have $s_n\ge(\eta/d)(1-o(1))N(\log
N)^\gamma$, and $c=\eta/(2d)$ works.
\end{proof}

\begin{theorem}[Theorem B]\label{thm:lower}
There are a language $\Wl\in\ENP$ and a constant $c>0$ such that for infinitely
many $N$ every circuit in $\cK(N^{3/4},N^{1/2})$ with $N$ inputs that decides
$\Wl$ on $\{0,1\}^N$ has more than $cN\sqrt{\log N}$ wires.
\end{theorem}

\begin{proof}
Theorem~\ref{thm:counting} with Remark~\ref{rem:space} decides satisfiability
by testing whether $\#C>0$ on a multitape machine in time
$(s+L)^{b_0+O(1)}2^{A_0s/\sqrt{\log s}}$. Apply Theorem~\ref{thm:transfer} with
$\gamma=\tfrac12$, and Lemma~\ref{lem:WinENP}.
\end{proof}

\begin{corollary}\label{cor:lowerclasses}
For infinitely many $N$, the language $\Wl$ on $\{0,1\}^N$ is not decided by any
circuit with fewer than $cN\sqrt{\log N}$ wires whose gates are gates of fan-in
two, AND, OR, parity, $\mathrm{MOD}_\mu$ and symmetric gates of any fan-in, and
threshold gates with weights below $2^{\sqrt N}$, in any combination and at any
depth.
\end{corollary}

\begin{proof}
Such circuits lie in $\cK(2,N^{1/2})\subseteq\cK(N^{3/4},N^{1/2})$ by
Example~\ref{ex:gates}.
\end{proof}

\subsection{Comparison}
For fan-in-two circuits, Theorem~B is \cite[Corollary 1.4]{OAI26}, since $g$
such gates use at most $2g$ wires. For threshold circuits, the bounds of
\cite{IPS97,KW16,CSS18} hold for functions in $\Pc$ and are stronger at each
fixed depth, but they tend to linear as the depth grows; Theorem~B is weaker at
small depth and holds at every depth. To the author's knowledge no superlinear
lower bound for a language in $\ENP$ against unbounded-depth circuits of
threshold or $\mathrm{MOD}_\mu$ gates, or of arbitrary gates of fan-in
$N^{3/4}$, was known before.

%%==================================================================
\section{Limits of the method}\label{sec:limits}
%%==================================================================

\subsection{The exponent $\sqrt{\log N}$}
Theorem~\ref{thm:transfer} converts a saving of $(\log s)^\gamma$ in the exponent
of a satisfiability algorithm into a factor $(\log N)^\gamma$ in the lower bound,
and conversely the method can give no more: the test circuit has $\Theta(s_n)$
wires and must be checked in time $2^{n/2}$. In the proof of
Theorem~\ref{thm:contraction} the factor $\sqrt{\log s}$ has a definite source
(Appendix~\ref{app:modes}). At a call with $M$ wires, an interpolation node of
two-adic level $i$ has about $M/D$ mode variables of each weight, $D=L/2^i$,
and works at a precision of order $MD/u$, where $u$ is the number of retained
levels; the number of admissible mode monomials is then
$\exp(O(\sqrt{M\cdot MD/u\cdot D^{-1}}))=\exp(O(M/\sqrt u))$. The number of
levels is at most $\log L$, and the packet length $L$ must stay below a small
power of $s$, because a packet tensor costs $\exp(O(L\log s))$ to evaluate.
So $u=O(\log s)$, and the recursion, which divides $M$ by eight at each depth,
sums these exponents to $O(s/\sqrt{\log s})$. A bound $cN\log N$ would need a contraction algorithm
whose exponent is $O(s/\log s)$, and the present design does not give one.

\subsection{Fan-in and weights}
Some restriction on the table gates is unavoidable, since one gate of fan-in
$N$ decides any language. In the proof, a table gate of fan-in $F$ costs $2^F$ in
the guess and in the entries of the network, and Theorem~\ref{thm:contraction}
tolerates entry costs of $2^{O(s^{9/10})}$; the exponents $4/5$, $2/3$ and $3/4$,
$1/2$ were chosen for simplicity and are not optimal. Weights enter through the
number $b$ of bits per block, which must keep the root degree below $s^{9/10}$.

\subsection{Superpolynomial bounds}
Williams' theorem \cite{Wil13} converts an algorithm for circuits of polynomial
size with $n$ inputs into a lower bound only if it beats $2^n$ by a
superpolynomial factor. Theorem~\ref{thm:counting} is measured in the size: for
a circuit with $n$ inputs and $n^k$ wires it gives time $2^{O(n^k/\sqrt{\log
n})}$, which is far above $2^n$. Superpolynomial lower bounds for $\ENP$, and a
fortiori for $\NP$, need algorithms whose running time depends on the number of
inputs; nothing in the tensor method provides this.

\subsection{From $\ENP$ to $\NP$}
The language $\Wl$ is defined by the first satisfying assignment of an
exponentially long formula and is computed with an $\NP$ oracle in exponential
time; it is not known to lie in $\NP$. The lower bound for
$\Wl$ therefore says nothing about $\Pc$ versus $\NP$ directly, and even a
superpolynomial lower bound for $\ENP$ would leave the question open, since
$\Pc\ne\NP$ is not implied by lower bounds for exponential-time classes.

\subsection{The barriers}
The proof diagonalises against $\NTIME(2^{n/2})$ through a single language, so
it is not a natural proof in the sense of \cite{RR97}. Because the table gates
may compute arbitrary functions of $N^{3/4}$ inputs, Theorem~B holds for circuits
with small gates of arbitrary function. The locality barrier \cite{CHOPRS22}
shows that the magnified statements of \cite{OS18} are false for circuits with
small oracle gates, so an argument that applies equally to such gates, as this
one does, is not expected to reach them. The barrier concerns techniques and
does not formally exclude any extension, but it indicates that the route from
Theorem~B to superpolynomial bounds does not pass through magnification.
